\documentclass[11pt,a4paper]{article}
\usepackage[margin=20mm]{geometry}
\usepackage[T1]{fontenc}
\usepackage{lmodern}
\usepackage{booktabs,tabularx,longtable,array,enumitem,hyperref,xcolor,listings}
\hypersetup{colorlinks=true,linkcolor=blue,urlcolor=blue}
\setlength{\parindent}{0pt}
\setlength{\parskip}{5pt}
\renewcommand{\arraystretch}{1.15}
\newcolumntype{Y}{>{\raggedright\arraybackslash}X}
\newcommand{\code}[1]{\texttt{\small #1}}
\begin{document}

\begin{titlepage}
\centering
{\Large 23CSE455 -- Design Patterns\par}
\vspace{12mm}
{\huge\bfseries Design Pattern Analysis in a\par SvelteKit-Based Smart Home Application\par}
\vspace{8mm}
{\Large Team A4\par}
\vspace{4mm}
{\large Case study: \textit{ha-fusion}\par}
\vfill
\begin{tabular}{ll}
\toprule
Roll number & Student name \\
\midrule
AM.AI.U4AID23027 & Adwaith S Menon (AIDS) \\
AM.SC.U4CSE23153 & Sai Samay (AIDS) \\
AM.AI.U4AID23044 & Karthikeya Sai (AIDS) \\
AM.AI.U4AID23028 & Amitha K Anilkumar (AIDS) \\
AM.SC.U4CSE23052 & Meenakshi Manoj (AIDS) \\
\bottomrule
\end{tabular}

\vspace{4mm}
\vfill
Submission repository: \url{https://github.com/aswathymohan-amrita/23CSE455_DP_CaseStudy}\\
Report date: 11 October 2026
\end{titlepage}

\tableofcontents
\newpage

\section{Application and architecture}
\subsection{Original source and environment}
The case study examines \textit{ha-fusion}, an open-source SvelteKit dashboard
for Home Assistant: \url{https://github.com/knowald/ha-fusion}. The frozen
baseline is commit \code{65cc1957c39c6005609102c33650c32abc8d12f0}; the recorded
license is MIT. The latest stable release tag at the frozen default-branch
HEAD is \code{2026.8.2}; its commit date is 2026-08-16. The inspection access
date is 2026-10-11. The submission
README records Node.js 22.15.0, pnpm 12.11.2, Svelte 5.56.8, SvelteKit 2.69.2,
Vite 8.1.4, and Vitest 4.1.10. Static analysis used cloc 2.10 and Lizard
1.24.1. The application uses Home Assistant as its external entity and control
service; the inspected project metadata does not identify a separate database
for the in-scope flows.

The baseline application source is preserved in \code{data/code/}; the evolved
source is in \code{Part2/code/}. Setup and test commands and inspected modules
are documented in \code{README.md}. The baseline test log reports 13 test files
and 46 tests passing. The final cumulative Part 2 logs report 15 test files and
68 tests passing, with \code{pnpm check} reporting zero errors and warnings.
These are saved run results; see \code{data/test\_result/} for the logs.

\subsection{Run and configuration commands}
Commands are run from the application root (\code{data/code/} for baseline,
\code{Part2/code/} for evolved source):
\begin{center}
\begin{tabularx}{\linewidth}{p{0.31\linewidth}p{0.18\linewidth}Y}
\toprule
Command & Recorded status & Result / note \\
\midrule
\code{pnpm install} & Run for baseline setup & Dependencies installed for the recorded baseline test run. \\
\code{pnpm dev -- --open} & Not run in this review & Requires a configured, reachable Home Assistant instance. \\
\code{pnpm build} & Not run in this review & No build result is claimed. \\
\code{pnpm test} & Baseline and evolved logs & Baseline: 13 files / 46 tests pass. Evolved final: 15 files / 68 tests pass. \\
\code{pnpm lint} & Not run in this review & No lint result is claimed. \\
\code{pnpm check} & Evolved final log & Zero errors and zero warnings in the saved final check. \\
\bottomrule
\end{tabularx}
\end{center}
The application README requires \code{HASS\_URL}; \code{TZ} is optional.
Development or end-to-end operation needs a running Home Assistant instance.
The baseline \code{package.json} does not declare a Node engine constraint;
the recorded environment uses Node.js 22.15.0 and pnpm 12.11.2.

\subsection{Inspection scope}
The pattern inventory covers the baseline collaboration records in
\code{data/patterns.csv}. Inspected areas include dashboard/entity stores,
recursive condition evaluation, Hearth configuration and history, entity modal
selection, Home Assistant subscriptions, and the Picture Elements canvas. The
full file/module inventory and exclusions (dependencies, generated output,
static assets, and unrelated tooling) are in \code{README.md} and the Part 1
evidence. Communication details are discussed only where needed to explain the
entity update flow; this report does not classify the protocol layer as a
separate design pattern.

\subsection{Architecture and representative flows}
The UI uses Svelte components and stores for dashboard and Hearth state.
Condition functions evaluate visibility trees; modal utilities select a
domain-specific entity modal. The Socket connection receives entity updates
from Home Assistant and publishes them through the \code{states} store for
reactive UI consumers.

\begin{lstlisting}[basicstyle=\ttfamily\small,frame=single]
User click
  -> Main/Button.svelte click handler and toggle()
  -> Utils.getTogglableService() chooses domain.service
  -> home-assistant-js-websocket.callService(connection, ...)
  -> Home Assistant service / device
  -> Socket.ts subscribeEntities callback
  -> Stores.states
  -> Main/Button.svelte $states effect updates entity view
\end{lstlisting}

Representative source traces recorded in the inventory are:
\begin{itemize}[nosep]
  \item Entity state updates: \code{Socket.ts:107--116} subscribes to entity
  changes; \code{Stores.ts:23} exposes the state store; the entity UI consumes
  that store (\code{Main/Entities.svelte:1--50}).
  \item Visibility: \code{Types.ts:58--69} defines the recursive condition
  shape; \code{Conditional.ts:70--90,132--163} dispatches and recursively
  evaluates conditions.
  \item Entity modal selection: \code{Hearth/modals.ts:11--171} chooses a
  modal using the entity domain, with domain extraction in \code{Utils.ts:74--76}
  and modal invocation in \code{Modals.ts:17--21}.
  \item Canvas behavior: \code{konvaBase.ts:13--517} supplies shared canvas
  lifecycle behavior specialized by \code{konvaViewer.ts} and
  \code{konvaEditor.ts}.
\end{itemize}

The device-control path is traceable in the baseline source: the selected
entity identifier is derived from \code{sel.entity\_id}
(\code{data/code/src/lib/Main/Button.svelte:41}); a non-edit-mode click invokes
\code{toggle()} (same file, lines 755--767 and 189--218). That function asks
\code{getTogglableService(entity)} to choose a domain-specific service
(\code{data/code/src/lib/Utils.ts:96--149}) and sends the selected service
with \code{entity\_id} through \code{callService(\$connection, ...)}
(\code{Main/Button.svelte:239--244}). The Socket connection subscribes to
Home Assistant entity updates and writes them to \code{states}
(\code{data/code/src/lib/Socket.ts:107--116}; \code{Stores.ts:21--24}); the
button effect reads \code{\$states[entity\_id]} to refresh its entity and
optimistic UI state (\code{Main/Button.svelte:100--118}). This traces the
application call to the Home Assistant library boundary; network delivery and
the remote device response are outside this source-level trace.

Preferences/configuration have a distinct persistence flow. The page server
loads \code{data/configuration.yaml} and \code{data/dashboard.yaml}, adds the
Home Assistant URL from the environment or proxy request, and returns them to
the UI (\code{data/code/src/routes/+page.server.ts:33--67}). Hearth edit state
is mutated through \code{updateConfig()} and saved by \code{saveEdit()} as a
revisioned POST (\code{data/code/src/lib/Hearth/store.ts:43--53,112--148});
the server serializes YAML and atomically writes the configuration file
(\code{save\_hearth/+server.ts:94--140}). These
configuration values affect dashboard layout, widgets, theme, localization,
and application settings. No separate database-backed preference store is
identified in this trace.

\section{Part 1 -- Pattern discovery and verification}
\subsection{Method and scope limits}
For Part 1, we used the frozen source to record participant roles, source
locations, and distinct collaborations. The saved historical prompt, model
output, and review are in \code{Part1/llm/}. We separated two independent
Memento collaborations that the initial output had combined and added source
locations to the inventory. A follow-up audit of the assigned smart-device,
control, and preference scope is now recorded in
\code{Part1/llm/p1\_discovery\_audit.md}; its confirmed-looking-only draft is
\code{Part1/llm/patterns\_p1\_discovery.csv}. The follow-up is a source-based
addition and is not presented as a historical model response.

That scoped audit records five confirmed-looking collaborations, two
uncertain candidates, and five rejected candidates. It is limited to the
inspected P1 paths listed in the audit note and makes no repository-wide
exhaustiveness claim. The canonical \code{data/patterns.csv} and its totals
remain the P2-owned baseline inventory; the P1 draft is a discovery handoff,
not a replacement for those final records.

\subsection{Baseline totals and inventory}
The recorded counts are five pattern types and eight pattern instances:
three Memento, two Composite, one Observer, one Factory Method, and one
Template Method. The count file records two distinct concrete participant
classes (\code{KonvaEditor} and \code{KonvaViewer}) under its stated counting
rule. A consolidated unique-participant-symbol total is \textbf{UNVERIFIED}:
the P2 inventory participant fields combine role aliases, variants, shorthand,
and grouped products and have not been normalized into a symbol-level ledger.
The scoped P1 draft separately deduplicates 32 participant symbols, as listed
in its audit note. The P2-owned canonical inventory contains only its confirmed
rows; the scoped P1 audit separately records two uncertain and five rejected
candidates.

\begin{longtable}{p{0.13\linewidth}p{0.25\linewidth}p{0.50\linewidth}}
\toprule
Pattern & Instance & Principal source participants and evidence location \\
\midrule
\endhead
Memento & Classic Dashboard Edit History & \code{dashboard}, \code{history}, \code{HistoryButtons.svelte}; \code{data/code/src/lib/Stores.ts:29,147--152}; \code{Drawer/HistoryButtons.svelte:18,43--82}. \\
Memento & Hearth Config History \& Edit Transaction & \code{hearthConfig}, snapshot stacks and edit operations; \code{data/code/src/lib/Hearth/store.ts:18,32--110}; configuration type in \code{Hearth/config.ts:324--341}. \\
Memento & Konva Editor History & \code{KonvaEditor} and \code{KonvaHistory}; \code{Modal/PictureElements/konvaEditor.ts:754--840}; type in \code{data/code/src/lib/Types.ts:350--353}. \\
Composite & Dashboard Section Layout Hierarchy & recursive \code{Section} and \code{findInSections}; \code{data/code/src/lib/Types.ts:42--56}; \code{data/code/src/lib/Utils.ts:25--60}. \\
Composite & Visibility Condition Tree & \code{Condition}, \code{validateCondition}, and \code{handleAnd}/\code{handleOr}; \code{data/code/src/lib/Types.ts:58--69}; \code{data/code/src/lib/Conditional.ts:58--216}. \\
Observer & Home Assistant WebSocket Entity Subscriptions & connection subscription, \code{states} store, and UI consumers; \code{data/code/src/lib/Socket.ts:107--116}; \code{data/code/src/lib/Stores.ts:23}; \code{data/code/src/lib/Main/Entities.svelte:1--50}. \\
Factory Method & Entity Domain Modal Factory & \code{openEntityModal}, \code{getDomain}, \code{openModal}, and concrete modal components; \code{data/code/src/lib/Hearth/modals.ts:11--171}; \code{Utils.ts:74--76}; \code{Modals.ts:17--21}. \\
Template Method & Picture Elements Konva Canvas Base & \code{KonvaBase}, \code{KonvaViewer}, and \code{KonvaEditor}; \code{konvaBase.ts:13--517}; \code{konvaViewer.ts} and \code{konvaEditor.ts} in the same directory. \\
\bottomrule
\end{longtable}

Full P2-owned role mappings are in \code{data/patterns.csv}; the scoped P1
discovery draft and disposition ledger are in \code{Part1/llm/}. Editable and
rendered architecture diagrams are in \code{Part1/pattern\_architecture/}.
The source audit note records that the first model output combined distinct
history mechanisms; the final inventory separates them. Strategy is not a
baseline type in this inventory; it first appears during Change 1.

\section{Part 2 -- Evolution and refactoring}
Three changes were implemented on the baseline and recorded in
\code{data/patterns\_evolution.csv}. Requirements, acceptance criteria, diffs,
and logs are under \code{Part2/code/changes/Change1--Change3/}.

\begin{longtable}{p{0.13\linewidth}p{0.30\linewidth}p{0.46\linewidth}}
\toprule
Change & Requirement and design work & Verification and measured effect \\
\midrule
\endhead
CH01 & Add recursive \code{not} evaluation and dispatch through a function-based Strategy map. Extends the Visibility Condition Tree Composite and introduces the evaluator Strategy. & Final cumulative suite: 68 tests pass; type check: zero errors/warnings. Measured \code{Conditional.ts} + \code{Types.ts}: 468 to 484 LOC; \code{validateCondition()} CCN 6 to 3. \\
CH02 & Extract bounded Hearth undo/redo state into generic \code{HistoryManager\textless T\textgreater} using Extract Class and Encapsulate Collection. & Final cumulative suite: 68 tests pass; clean type check. Measured Hearth directory: 17,063 to 17,186 LOC; maximum history-operation CCN remains 2. \\
CH03 & Add configurable alert rail widget and dispatch configuration predicates through \code{WIDGET\_CONFIG\_VALIDATORS}. & Final cumulative suite: 68 tests pass; clean type check. Measured Hearth directory: 17,186 to 17,293 LOC; \code{railWidgetNeedsConfiguration()} CCN 7 to 2. \\
\bottomrule
\end{longtable}

Pattern effects, source-version totals, limitations, and exact logs are recorded
in \code{Part2/Part2\_Report.md}, \code{data/patterns\_evolution.csv},
\code{data/pattern\_counts.csv}, and \code{Part2/evolution\_metrics.csv}.
Across tracked snapshots, the inventory moves from 5 types/8 instances at the
baseline to 6 types/10 instances after CH03. One CH03 Factory Method candidate
was rejected because the registry and renderer branches do not show the
Factory Method collaboration. Baseline pattern rows were carried forward in
the evolution table; CH01--CH03 did not perform a new exhaustive Part 1 sweep.

\begin{center}
\begin{tabular}{lrrrrr}
\toprule
Snapshot & Types & Instances & Concrete classes & Rejected & Uncertain \\
\midrule
Original & 5 & 8 & 2 & 0 & 0 \\
After CH01 & 6 & 9 & 2 & 0 & 0 \\
After CH02 & 6 & 9 & 3 & 0 & 0 \\
After CH03 & 6 & 10 & 3 & 1 & 0 \\
\bottomrule
\end{tabular}
\end{center}
These are tracked-inventory counts. Baseline rejected/uncertain zeros mean no
such rows were carried into the inventory; they do not represent the result
of an exhaustive baseline candidate sweep.

\section{Part 3 -- Cross-language comparison}
Six selected collaborations were implemented as small behavior-matched
examples in Java, Python, JavaScript, and C++: Composite / Visibility Condition
Tree; Observer / Entity Subscriptions; Memento / Hearth History; Strategy /
Visibility Condition Evaluator; Factory Method / Entity Modal; and Template
Method / Canvas Renderer. The common fixtures and expected behavior are in
\code{Part3/test-cases/common\_spec.md}; production snippets and separate
runners are grouped by pattern and language.

\begin{center}
\begin{tabular}{lrrrrl}
\toprule
Language & Mean LOC & Mean max CCN & Saved suites & Local rerun & Mean saved ms \\
\midrule
Java & 68.0 & 5.0 & 6/6 pass & 6/6 pass & 8.196 \\
Python & 59.7 & 4.8 & 6/6 pass & 6/6 pass & 0.028 \\
JavaScript & 61.3 & 3.7 & 6/6 pass & 6/6 pass & 0.301 \\
C++ & 97.0 & 4.7 & 6/6 pass & not rerun here & 0.038 \\
\bottomrule
\end{tabular}
\end{center}

The canonical per-pattern/per-language data, exact runtimes and commands are
in \code{data/cross\_language.csv}; test logs and recomputed static metrics
are in \code{Part3/evidence/}. Python, JavaScript, and Java runners passed
locally (18 suites) on Python 3.12.2, Node.js 22.15.0, and OpenJDK 17.0.10.
The six C++ runners were not rerun on this macOS host because their timing
harnesses require Windows headers and performance-counter APIs; saved Windows
logs report their earlier MinGW runs passing.

LOC excludes blanks and comments. Maximum function CCN is measured with Lizard
1.24.1; C++ counts include both implementation and header. The saved timing
values are medians of five measured runner-body runs after one warm-up. They
exclude process creation and VM/interpreter startup. The original timing
hardware and operating system were not preserved, so these small test timings
are not a controlled cross-language performance ranking. The earlier slide
PDF is retained as historical output; its C++ LOC omitted headers.

Java uses declared interfaces/classes and runtime dispatch; Python supports
abstract base classes and dynamic/structural use; JavaScript class syntax uses
prototype-based objects and dynamic dispatch; C++ examples use abstract base
classes and virtual functions. These are language-specific expression choices,
not evidence that one language is superior. References to the relevant
language documentation are listed at the end.

\section{LLM records and human verification}
AI-assisted work and its verification are recorded in \code{data/llm.csv};
supporting prompts, inputs, and review notes are stored with each part. The
Part 1 record describes how we separated the merged Memento collaborations
and added source locations. The Part 2 completion record notes that the
historical CH01--CH03 model prompts and outputs were unavailable. The current
reference task specification is labeled as reconstructed, not historical. For Part 3, the
earlier raw model response was unavailable. Its prompt/context and the later
completion review are preserved, and the later review is explicitly not
presented as the earlier raw response. Evidence and limitations are in
\code{Part2/llm/verification.md} and \code{Part3/evidence/verification.md}.

The current Part 1 task specification is stored in
\code{Part1/Part1\_prompt.txt}. The shorter historical Gemini prompt is
preserved separately at \code{Part1/llm/historical\_discovery\_prompt.txt}
and is the prompt linked to the saved historical output. Part 2's reference
prompt is explicitly marked as reconstructed; it is not presented as a recovered
historical transcript.

\section{Findings, contributions, and limitations}
The baseline inventory describes undo/redo in three separate subsystems,
recursive compositions for layouts and conditions, a domain-specific modal
selection collaboration, an entity-state subscription flow, and a shared
canvas lifecycle specialized by viewer/editor subclasses. The evolution added
condition negation, extracted history management, and introduced an alert
widget; test evidence records behavior preservation and new behavior. The
cross-language examples retain the same stated fixtures, but their small test
workloads and differing historical runtime environments limit conclusions
about performance.

\begin{longtable}{p{0.27\linewidth}p{0.66\linewidth}}
\toprule
Team member & Contribution recorded from the team's work-division note \\
\midrule
\endhead
Adwaith S Menon & Part 1 pattern discovery. \\
Sai Samay & Part 1 pattern discovery (shared with Adwaith S Menon). \\
Karthikeya Sai & Part 1 pattern /inventory; Part 2 evolution and refactoring. \\
Amitha K Anilkumar & Part 3 implementations in Java, Python, and JavaScript. \\
Meenakshi Manoj & Part 3 C++ implementation and report writing. \\
\bottomrule
\end{longtable}

The team confirmed that Amitha K Anilkumar implemented the Part 3 examples in
Java, Python, and JavaScript, while Meenakshi Manoj handled the report and C++
implementation. Karthikeya Sai handled Part 1 verification and Part 2 evolution/refactoring, and all five members
are in the AIDS batch. The other contribution descriptions follow the work-
division details supplied by the team.

\subsection{References}
\begin{itemize}[nosep]
\item Ha-Fusion repository and baseline: \url{https://github.com/knowald/ha-fusion}; MIT License.
\item Oracle Java Tutorials, Polymorphism: \url{https://docs.oracle.com/javase/tutorial/java/IandI/polymorphism.html}.
\item Oracle Java Tutorials, Defining an Interface: \url{https://docs.oracle.com/javase/tutorial/java/IandI/interfaceDef.html}.
\item Python 3.12 documentation, Abstract Base Classes: \url{https://docs.python.org/3.12/library/abc.html}.
\item MDN Web Docs, JavaScript Classes: \url{https://developer.mozilla.org/en-US/docs/Web/JavaScript/Reference/Classes}.
\item C++ working draft, virtual functions: \url{https://eel.is/c++draft/class.virtual}.
\item Gamma, Helm, Johnson, and Vlissides, \textit{Design Patterns: Elements of Reusable Object-Oriented Software}, Addison-Wesley, 1994.
\item Larman, \textit{Applying UML and Patterns}, 3rd ed., Prentice Hall, 2004.
\item cloc 2.10 and Lizard 1.24.1 measurement tools; versions and rules are documented in \code{README.md} and \code{Part3/evidence/measurement\_notes.md}.
\end{itemize}

\appendix
\section{Submission contents and checks}
This team folder contains the consolidated report source and compiled PDF,
README, the canonical datasets in \code{data/}, original source and test logs,
Part 1 diagrams and evidence, Part 2 changes and snapshots, and Part 3
snippets, shared test cases, evidence, and report. The course instructions request one \code{A4/} team folder and one pull
request to the upstream repository's main branch. Capture the requested repository screenshot after
the team folder is uploaded. Before submission, make sure the PDF is compiled from the final version of
this TeX source.

\end{document}
